% Standalone problem document, generated from the adjacent README.md.
% Compile with XeLaTeX. Mathematical content is maintained in Markdown.
\documentclass[11pt,a4paper]{article}
\usepackage[fontsize=10.5pt]{scrextend}
\usepackage[margin=22mm,headheight=15pt,headsep=9mm,footskip=11mm]{geometry}
\usepackage{fontspec}
\setmainfont{texgyrepagella}[Extension=.otf,UprightFont=*-regular,BoldFont=*-bold,ItalicFont=*-italic,BoldItalicFont=*-bolditalic]
\setsansfont{texgyreheros}[Extension=.otf,UprightFont=*-regular,BoldFont=*-bold,ItalicFont=*-italic,BoldItalicFont=*-bolditalic]
\setmonofont{texgyrecursor}[Extension=.otf,UprightFont=*-regular,BoldFont=*-bold,ItalicFont=*-italic,BoldItalicFont=*-bolditalic,Scale=MatchLowercase]
\usepackage{amsmath,amssymb}
\usepackage{unicode-math}
\setmathfont{texgyrepagella-math.otf}
\usepackage{xcolor,microtype,enumitem,needspace,fancyhdr}
\usepackage{bookmark}
\definecolor{ink}{HTML}{17344B}
\definecolor{accent}{HTML}{197D80}
\definecolor{muted}{HTML}{596773}
\hypersetup{colorlinks=true,linkcolor=accent,urlcolor=accent,pdftitle={NM-01 - Polynomial-time
minimum-volume decision under sufficient
scattering},pdfauthor={Open Problems in Numerical Linear Algebra}}
\urlstyle{same}
\setlength{\parindent}{0pt}
\setlength{\parskip}{5pt plus 1pt minus 1pt}
\linespread{1.04}
\setlength{\emergencystretch}{3em}
\widowpenalty=10000
\clubpenalty=10000
\displaywidowpenalty=10000
\allowdisplaybreaks[1]
\setlist{nosep,leftmargin=1.5em}
\providecommand{\tightlist}{\setlength{\itemsep}{0pt}\setlength{\parskip}{0pt}}
\providecommand{\pandocbounded}[1]{#1}
\pagestyle{fancy}
\fancyhf{}
\fancyhead[L]{\sffamily\footnotesize\color{muted}OPEN PROBLEMS IN NUMERICAL LINEAR ALGEBRA}
\fancyhead[R]{\sffamily\bfseries\footnotesize\color{accent}NM-01}
\fancyfoot[L]{\parbox[b]{0.9\textwidth}{\sffamily\fontsize{8}{10}\selectfont\color{muted}Editorial ratings; a bounded search cannot certify that a problem remains unsolved.\\Literature check: 2026-09-08}}
\fancyfoot[R]{\sffamily\footnotesize\color{muted}\thepage}
\renewcommand{\headrulewidth}{0.3pt}
\renewcommand{\headrule}{\hbox to\headwidth{\color{accent}\leaders\hrule height \headrulewidth\hfill}}
\makeatletter
\renewcommand{\section}{\Needspace{5\baselineskip}\@startsection{section}{1}{0pt}{12pt plus 2pt minus 2pt}{4pt}{\sffamily\bfseries\large\color{ink}}}
\renewcommand{\subsection}{\Needspace{4\baselineskip}\@startsection{subsection}{2}{0pt}{9pt plus 1pt minus 1pt}{3pt}{\sffamily\bfseries\normalsize\color{ink}}}
\makeatother
\setcounter{secnumdepth}{0}
\begin{document}
{\sffamily\small\color{muted}Nonnegative and positive
factorizations\par}
\vspace{3pt}
{\raggedright\hyphenpenalty=10000\exhyphenpenalty=10000\sffamily\bfseries\fontsize{21}{24}\selectfont\color{ink}Polynomial-time
minimum-volume decision under sufficient scattering\par}
\vspace{7pt}
{\sffamily\small\textbf{Difficulty:} extreme\quad\textbf{Importance:} broadly
interesting\par}
\vspace{4pt}
{\color{accent}\hrule height 0.8pt}
\vspace{6pt}
\textbf{Topic:} structured matrix factorization; global optimization

\textbf{Status:} open; admitted after a bounded literature search found
no resolution.

\subsection{Context and notation}\label{context-and-notation}

For \(H\in\mathbb R_+^{r\times n}\), let
\(\operatorname{cone}(H)=\{Hy:y\in\mathbb R_+^n\}\) and let \(e_d\) be
the all-ones vector in \(\mathbb R^d\). Here \textbf{SSC} means the
precise version in Gillis's Definition 4.15:

\[
\mathcal C_r=\{x\in\mathbb R_+^r:
e_r^Tx\geq\sqrt{r-1}\|x\|_2\}\subseteq\operatorname{cone}(H),
\]

and every orthogonal \(Q\in\mathbb R^{r\times r}\) satisfying
\(\operatorname{cone}(H)\subseteq\operatorname{cone}(Q)\) is a
permutation matrix. Some later sources use a stronger second condition
involving the boundary of the dual cone; that condition is not
substituted here.

For \(X\in\mathbb R^{m\times n}\) and \(r=\operatorname{rank}(X)\), the
minimum-volume problem considered here is

\[
\min_{W\in\mathbb R^{m\times r},\ H\in\mathbb R_+^{r\times n}}
\det(W^TW)
\quad\text{subject to}\quad
X=WH,\qquad e_r^TH=e_n^T.\tag{MV}
\]

In particular, \(W\) is unrestricted in sign and the \textbf{columns} of
\(H\) sum to one. This is min-vol NMF (1), Definition 4.42 of the book.

\subsection{Problem statement}\label{problem-statement}

Consider the following rational-input decision version of minimum-volume
optimization. The input is \(X\in\mathbb Q^{m\times n}\), an integer
\(2\leq r\leq\min(m,n)\), and \(\tau\in\mathbb Q_{\geq0}\), with the
promise that \(\operatorname{rank}(X)=r\) and that there exist real
factors \(X=W_\star H_\star\) feasible for (MV) with \(H_\star\)
satisfying SSC. These factors are not supplied. Decide whether

\[
\exists W\in\mathbb R^{m\times r},\ H\in\mathbb R_+^{r\times n}:
\quad X=WH,\quad e_r^TH=e_n^T,\quad\det(W^TW)\leq\tau.
\]

Does a polynomial-time algorithm exist, with time measured in the total
binary input length? Randomization is allowed with success probability
at least \(2/3\) on every promised input. The algorithm must run in
polynomial time on all inputs but need only answer correctly on promised
inputs. Checking the SSC promise is not part of the task. This decision
formulation fixes the computational model of the book's
polynomial-solvability conjecture; it does not assert an equivalence
with exact factor recovery.

\newpage
\subsection{References}\label{references}

Gillis, \href{https://doi.org/10.1137/1.9781611976410}{\emph{Nonnegative
Matrix Factorization}}, §4.3.3.6, pp.~148--149, with Definitions 4.15
and 4.42 and Theorem 4.43;
\href{https://orbi.umons.ac.be/bitstream/20.500.12907/42337/1/NMFbook_SIAM_reprint.pdf}{author-hosted
book}. Barbarino, Gillis, and Saha,
\href{https://arxiv.org/html/2511.04291v1}{\emph{Robustness of
Minimum-Volume Nonnegative Matrix Factorization under an Expanded
Sufficiently Scattered Condition}}, §5, final research question.

\subsection{Status check}\label{status-check}

Searches for
\textit{minimum-volume NMF polynomial time sufficiently scattered}
and
\textit{minimum volume simplex sufficiently scattered algorithm 2026 2025}
found no polynomial-time guarantee for this promise. The November 2025
paper still asks for complexity results under its stronger \(p\)-SSC
assumption. Its robustness theorems assume a globally optimal
minimum-volume solution; they do not compute one in polynomial time. The
book explains why the maximum-inscribed-ellipsoid approach can require
exponentially many polytope facets.
\end{document}
